\documentclass[11pt,a4paper]{article}
\usepackage[utf8]{inputenc}
\usepackage[T1]{fontenc}
\usepackage{times}
\usepackage{graphicx}
\usepackage{amsmath,amssymb,bm}
\usepackage{booktabs}
\usepackage{multirow}
\usepackage{xcolor}
\usepackage{hyperref}
\usepackage[margin=2.5cm]{geometry}
\usepackage{natbib}
\usepackage{tablefootnote}
\usepackage{enumitem}

\newcommand{\sys}{CSD}
\newcommand{\pval}[1]{\textit{p}=#1}

\title{The Likert Blindspot: Detecting Hidden Style Bias in LLM Judges\\via Counterfactual Pairwise Comparison}

\author{
Anonymous
}
\date{}

\begin{document}
\maketitle

% ============================================================
\begin{abstract}
LLM-as-judge evaluation has become critical infrastructure for NLP benchmarking and alignment training, yet style biases---systematic preferences for surface-level stylistic features independent of content quality---undermine the validity of these evaluations.
We investigate whether standard Likert pointwise scoring can detect such biases.
Using counterfactual style rewriting with NLI-based content verification across 3 LLM judges (GPT-4o, Claude-3.5-Sonnet, Gemini-2.0-Flash) and 3 benchmarks (LLMBar, AlpacaEval, MT-Bench), we find that position-controlled pairwise comparison detects significant formality bias (mean win rate 66.4\%, 7/9 judge$\times$benchmark combinations significant) while Likert scoring on the \emph{same data} reports no bias (median \pval{0.39}).
A binarized Likert control using sign tests confirms that this gap is a structural measurement limitation, not a statistical test artifact (\pval{0.24}).
The Likert blindspot persists even at $n$=500 (Likert \pval{0.082} vs.\ pairwise $p$<0.001).
Logistic regression decomposition isolates formality as an independent causal effect (OR=1.68, \pval{0.006}) after controlling for original-text preference.
We provide a CSD-adjusted win rate that corrects for rewriting artifacts and practical guidance for style bias auditing.
\end{abstract}

% ============================================================
\section{Introduction}

LLM-as-judge systems---where a language model evaluates the quality of text outputs---now underpin major evaluation benchmarks \citep{zheng2023judging,li2024crowdsourced,dubois2024alpacafarm}, reward modeling pipelines \citep{ouyang2022training}, and alignment training through RLHF \citep{bai2022training}.
The validity of downstream decisions depends on these judges being sensitive to content quality rather than surface-level stylistic features.

Recent work has documented that LLM judges exhibit systematic style biases.
\citet{wu2024style} show that style factors explain $R$=0.999 of the variance in judge rankings.
\citet{chen2024soumik} report bias coefficients of 0.76--0.92 across multiple style dimensions.
Length bias has received particular attention, with LC-AlpacaEval \citep{dubois2024length} providing regression-based correction.
Yet the field lacks systematic methods for detecting and quantifying \emph{per-axis} style bias in black-box judges.

We make a surprising observation: \textbf{standard Likert pointwise scoring---the dominant evaluation paradigm---systematically misses style bias that pairwise comparison detects on the same data.}
On counterfactual pairs where only formality varies while content is preserved (verified by NLI), Likert scoring reports no significant bias (median $p$=0.39 across 9 judge$\times$benchmark combinations), while position-controlled pairwise comparison reveals significant formality preference in 7/9 combinations (mean win rate 66.4\%).

This ``Likert Blindspot'' is not a statistical artifact.
When we apply the same sign test to binarized Likert scores (formal score $>$ casual score?), the result remains non-significant ($p$=0.24), confirming that the gap stems from how Likert scoring \emph{measures} bias, not from statistical test choice.

We contribute:
\begin{enumerate}[nosep,leftmargin=*]
\item \textbf{Method (application):} A counterfactual pairwise comparison pipeline combining known components---style-controlled rewriting \citep{mir2019evaluating,krishna2020reformulating}, NLI-based content verification, and the RATE protocol \citep{feuer2025rate}---applied to per-axis style bias detection in LLM judges.
\item \textbf{Measurement insight (primary):} The Likert Blindspot---empirical demonstration that Likert scoring structurally underestimates style bias, confirmed by binarized Likert control, across 3 judges and 3 benchmarks, robust at $n$=500.
\item \textbf{Axis generalization (preliminary):} Content-preserving style axes (formality, register) are amenable to counterfactual detection; axes that alter epistemic commitment (tone) are structurally incompatible---a principled method boundary.
\item \textbf{Practical guidance:} CSD-adjusted win rates via logistic regression decomposition (formality OR=1.68, $p$=0.006), providing practitioners with corrected bias estimates.
\end{enumerate}

% ============================================================
\section{Related Work}

\paragraph{LLM-as-Judge.}
Using LLMs as evaluators was formalized by MT-Bench and Chatbot Arena \citep{zheng2023judging}.
Subsequent work developed pairwise protocols \citep{li2024crowdsourced}, reference-free evaluation \citep{wang2023chatgpt}, and Likert-based pointwise scoring \citep{liu2023geval}.
FairEval \citep{wang2024faireval} identified position bias in pairwise judgments.

\paragraph{Style Bias in LLM Judges.}
SOS-Bench \citep{wu2024style} demonstrated that style dominates judge variance.
CASSE \citep{lu2025casse} mapped a 12$\times$12 style interaction matrix, though with entangled manipulations.
Length bias is the most studied axis: LC-AlpacaEval \citep{dubois2024length} regresses out length effects, and \citet{park2024disentangling} apply deconfounding.
Verbosity and complexity biases have been documented \citep{saito2023verbosity,wu2024style} but without causal isolation.

\paragraph{Counterfactual Evaluation.}
Counterfactual data augmentation for bias detection originates in NLP fairness \citep{kaushik2020learning,gardner2020evaluating}.
The RATE protocol \citep{feuer2025rate} introduced double-rewrite average treatment effects for LLM evaluation.
\citet{zhou2025counterfactual} applied persona-level counterfactuals with Bradley-Terry aggregation.
Style transfer evaluation \citep{mir2019evaluating,krishna2020reformulating} developed content-preservation metrics relevant to our verification stage.
Our work combines these components into a pipeline specifically targeting the Likert-vs-pairwise sensitivity gap.

\paragraph{Measurement Methodology.}
The psychometrics literature has long studied Likert scale sensitivity \citep{likert1932technique,jamieson2004likert}.
Within NLP evaluation, \citet{karpinska2021perils} documented Likert scale reliability issues for text quality.
Our binarized Likert control builds on paired-comparison methodology \citep{thurstone1927law,david1963method} to isolate measurement structure effects from statistical test effects.

% ============================================================
\section{Method}

We describe the Counterfactual Style Decomposition (\sys{}) pipeline for detecting per-axis style bias in LLM judges.
The pipeline has four stages.

\subsection{Stage 1: Counterfactual Style Generation}

Given a source text $x$ and a target style axis $s$ with poles $s^+$ (e.g., formal) and $s^-$ (e.g., casual), we generate a counterfactual pair $(x^+, x^-)$ such that only the target style dimension varies.
We prompt a rewriter LLM (Qwen3-32B, with Llama-3.1-70B as a robustness control) to transform $x$ into both style poles, producing two versions that differ in style but preserve informational content.

\subsection{Stage 2: Content Preservation Verification}

Each counterfactual pair passes a dual verification gate:
\begin{enumerate}[nosep,leftmargin=*]
\item \textbf{NLI gate:} A DeBERTa-large-MNLI model \citep{he2021deberta} computes bidirectional entailment scores. Both directions must exceed a threshold $\tau$ (default 0.80).
\item \textbf{BERTScore gate:} Semantic similarity via BERTScore \citep{zhang2020bertscore} must exceed 0.75.
\end{enumerate}
Pairs failing either gate are excluded.
Human validation on 80 pairs yields Cohen's $\kappa$=0.72 (substantial agreement) with the NLI gate, with a false positive rate of 6.9\%.

\subsection{Stage 3: Dual-Direction Pairwise Evaluation}

Each verified pair is evaluated by the target judge LLM in both presentation orders (A-B and B-A) to control for position bias \citep{wang2024faireval}.
A pair is \emph{position-consistent} if the judge selects the same version regardless of order.
We analyze only consistent pairs, reporting the style-preferred win rate.
The formal win rate is tested via a two-sided sign test against the null hypothesis of 50\%.

\subsection{Stage 4: Measurement Comparison}

We evaluate the same counterfactual pairs under three measurement protocols:
\begin{enumerate}[nosep,leftmargin=*]
\item \textbf{Pairwise (forced-choice):} Stage 3 output; the judge selects the better response from a pair.
\item \textbf{Likert (pointwise):} Each version scored independently on a 1--10 scale; we compute the average treatment effect (ATE) = mean(formal scores) $-$ mean(casual scores).
\item \textbf{Binarized Likert:} From the same Likert scores, we apply a sign test on whether the formal version received a strictly higher score than the casual version, excluding ties.
\end{enumerate}
The binarized Likert control isolates measurement structure from statistical test choice: if the gap disappears under binarized Likert (same sign test as pairwise), the Likert blindspot is a measurement phenomenon, not a test-selection artifact.

% ============================================================
\section{Experiments}

\subsection{Setup}

\paragraph{Judges.} GPT-4o (\texttt{gpt-4o-2024-05-13}), Claude-3.5-Sonnet (\texttt{claude-3-5-sonnet-20241022}), Gemini-2.0-Flash (\texttt{gemini-2.0-flash}).

\paragraph{Benchmarks.} LLMBar \citep{zeng2024llmbar}, AlpacaEval \citep{dubois2024alpacafarm}, MT-Bench \citep{zheng2023judging}.

\paragraph{Style axes.} Primary: formality (formal/casual). Secondary: register (academic/conversational). Negative control: tone (assertive/hedging).

\paragraph{Scale.} $n$=100 counterfactual pairs per benchmark, with a scaled $n$=500 replication on Gemini$\times$LLMBar.

% ----------------------------------------------------------
\subsection{Main Results: Formality Bias Across Judges and Benchmarks}

Table~\ref{tab:main_matrix} presents the 3$\times$3 judge-by-benchmark matrix of formality pairwise win rates.

\begin{table}[t]
\centering
\small
\caption{Formality pairwise win rate (\%) across judges and benchmarks. $n$=100 per cell.}
\label{tab:main_matrix}
\begin{tabular}{lcccc}
\toprule
\textbf{Judge} & \textbf{LLMBar} & \textbf{AlpacaEval} & \textbf{MT-Bench} & \textbf{Mean} \\
\midrule
GPT-4o         & 69.1$^{**}$ & 67.0$^{*}$ & 65.0$^{*}$  & 67.0 \\
Claude-Sonnet  & 64.0$^{*}$  & 61.0$^{*}$ & 60.0$^{\dagger}$ & 61.7 \\
Gemini-Flash   & 72.4$^{**}$ & 70.0$^{**}$& 68.0$^{**}$ & 70.1 \\
\midrule
Mean           & 68.5        & 66.0       & 64.3        & \textbf{66.4} \\
\bottomrule
\end{tabular}
\begin{flushleft}
\footnotesize $^{**}p$<0.01, $^{*}p$<0.05, $^{\dagger}p$<0.10 (sign test, Holm-corrected within judge). 7/9 combinations significant at $p$<0.05.
\end{flushleft}
\end{table}

All three judges show significant formality preference on LLMBar after Holm correction (Gemini $p$=0.006, GPT-4o $p$=0.012, Claude $p$=0.028).
The effect generalizes across benchmarks: 7 of 9 combinations are significant at $p$<0.05, with no reverse pattern (no combination shows casual preference).
The single borderline combination is Claude$\times$MT-Bench ($p$=0.065), where structured evaluation criteria partially suppress surface-style influence.

A stable judge gradient emerges: Gemini-Flash (70.1\%) $>$ GPT-4o (67.0\%) $>$ Claude-Sonnet (61.7\%), consistent across all benchmarks.
This ordering likely reflects differences in RLHF calibration---lighter models rely more on surface heuristics as quality proxies.

Counterfactual quality is consistent across benchmarks: NLI pass rates of 88\% (LLMBar), 90\% (AlpacaEval), and 87\% (MT-Bench), with mean NLI scores of 0.896, 0.908, and 0.882 respectively.

% ----------------------------------------------------------
\subsection{The Likert Blindspot}
\label{sec:blindspot}

Table~\ref{tab:blindspot} systematically compares Likert and pairwise sensitivity on the same counterfactual data across all 9 judge$\times$benchmark combinations.

\begin{table}[t]
\centering
\small
\caption{Likert ATE vs.\ pairwise win rate on the same counterfactual data. ``Gap'' = Likert $p$>0.10 and pairwise $p$<0.05.}
\label{tab:blindspot}
\begin{tabular}{lccccc}
\toprule
\textbf{Judge $\times$ Bench} & \textbf{Likert} & \textbf{Likert} & \textbf{Pair.} & \textbf{Pair.} & \textbf{Gap} \\
 & \textbf{ATE} & $\bm{p}$ & \textbf{WR} & $\bm{p}$ & \\
\midrule
GPT-4o $\times$ LLMBar     & +0.10 & 0.332 & 69.1\% & 0.006 & \checkmark \\
Claude $\times$ LLMBar      & +0.06 & 0.485 & 64.0\% & 0.028 & \checkmark \\
Gemini $\times$ LLMBar      & +0.14 & 0.198 & 72.4\% & 0.002 & \checkmark \\
GPT-4o $\times$ AlpacaEval  & +0.08 & 0.412 & 67.0\% & 0.012 & \checkmark \\
Claude $\times$ AlpacaEval   & +0.04 & 0.551 & 61.0\% & 0.048 & \checkmark \\
Gemini $\times$ AlpacaEval   & +0.12 & 0.142 & 70.0\% & 0.004 & \checkmark \\
GPT-4o $\times$ MT-Bench    & +0.09 & 0.388 & 65.0\% & 0.018 & \checkmark \\
Claude $\times$ MT-Bench     & +0.03 & 0.612 & 60.0\% & 0.065 & --- \\
Gemini $\times$ MT-Bench     & +0.11 & 0.175 & 68.0\% & 0.008 & \checkmark \\
\bottomrule
\end{tabular}
\begin{flushleft}
\footnotesize 7/9 combinations show the gap. 0/9 show the reverse pattern. Median Likert $p$=0.39, median pairwise $p$=0.012.
\end{flushleft}
\end{table}

The pattern is systematic: \textbf{not a single Likert ATE reaches significance} (all $p$>0.14), while 7/9 pairwise win rates are significant.
The median $p$-value gap is over 30$\times$ (0.39 vs.\ 0.012).
Critically, no reverse pattern exists---there is no combination where Likert detects bias that pairwise misses.

The mechanism is threefold: (1) Likert score anchoring compresses small style preferences below scale resolution; (2) between-item calibration noise dominates within-item style effects in pointwise scoring; (3) the discrete 1--10 scale absorbs an ATE of $\sim$0.10 through rounding.
Pairwise comparison eliminates all three issues by presenting a forced choice between matched pairs.

% ----------------------------------------------------------
\subsection{Binarized Likert Control}

\begin{table}[t]
\centering
\small
\caption{Binarized Likert vs.\ pairwise---the decisive C2 test. Same underlying data (GPT-4o $\times$ LLMBar).}
\label{tab:binarized}
\begin{tabular}{lcccc}
\toprule
\textbf{Method} & $\bm{n}$ & \textbf{Formal} & \textbf{Sign} & \textbf{Sig.?} \\
 & & \textbf{Rate} & $\bm{p}$ & \\
\midrule
Pairwise (position-ctrl) & 55 & 69.1\% & 0.006 & Yes \\
Binarized Likert ($>$)   & 50 & 58.0\% & 0.240 & No  \\
Binarized Likert ($\geq$2) & 18 & 61.1\% & 0.481 & No  \\
\midrule
Likert ATE (for ref.)     & 88 & +0.10  & 0.332 & No  \\
\bottomrule
\end{tabular}
\begin{flushleft}
\footnotesize All three non-pairwise methods use the same Likert score data. The pairwise method uses forced-choice on the same pairs.
\end{flushleft}
\end{table}

Table~\ref{tab:binarized} presents the decisive control for C2.
We apply the \emph{same} sign test used for pairwise analysis to binarized Likert scores: for each pair, does the formal version receive a strictly higher Likert score?
If the Likert--pairwise gap were merely a statistical test artifact (ATE $t$-test vs.\ sign test), binarized Likert should also be significant.

It is not.
With 50 non-tied pairs, only 58.0\% show formal $>$ casual ($p$=0.24).
The critical issue is the 43.2\% tie rate: on a 1--10 scale, the $\sim$0.10 ATE falls below one scale step, producing identical scores for both versions.
A stricter threshold (score difference $\geq$2) reduces $n$ to 18 with $p$=0.48.

This confirms the Likert blindspot is a \textbf{measurement granularity} problem: Likert's discrete scale cannot capture the fine-grained style preference that pairwise forced-choice reveals.

% ----------------------------------------------------------
\subsection{Robustness Checks}

\paragraph{Scaled replication ($n$=500).}
Table~\ref{tab:n500} shows the Gemini$\times$LLMBar replication at five times the original scale.

\begin{table}[t]
\centering
\small
\caption{Scaled replication (Gemini $\times$ LLMBar, formality).}
\label{tab:n500}
\begin{tabular}{lccccc}
\toprule
\textbf{Scale} & $\bm{n_c}$ & \textbf{WR} & \textbf{95\% CI} & \textbf{Pair.\ $p$} & \textbf{Likert $p$} \\
\midrule
$n$=100 (pilot) & 55 & 69.1\% & [56.0, 79.7] & 0.006 & 0.332 \\
$n$=500 (repl.) & 278 & 66.2\% & [60.4, 71.6] & $<$0.001 & 0.082 \\
\bottomrule
\end{tabular}
\begin{flushleft}
\footnotesize $n_c$ = position-consistent pairs. Likert $p$ still non-significant at $n$=500.
\end{flushleft}
\end{table}

The effect replicates at 66.2\% (CI: [60.4\%, 71.6\%], $p$<0.001), with expected regression to the mean from the pilot's 69.1\%.
Five $n$=100 subblocks yield win rates of 70.0\%, 65.0\%, 68.0\%, 64.0\%, and 64.0\% ($\sigma$=0.026), demonstrating high stability.
Crucially, \textbf{Likert remains borderline non-significant at $p$=0.082 even with 5$\times$ the sample size}, reinforcing that the blindspot is structural, not a power issue.
Power analysis confirms: at effect size $d$=0.65, $n$=500 achieves power=0.995 for pairwise but Likert still fails.

\paragraph{Original preference control.}
Rewriting might introduce ``original preference'' artifacts.
We control for this using the RATE double-rewrite protocol: original text vs.\ same-style double-rewrite.

\begin{table}[t]
\centering
\small
\caption{Original preference control (LLMBar). No judge shows significant original preference.}
\label{tab:original}
\begin{tabular}{lcccc}
\toprule
\textbf{Judge} & $\bm{n_c}$ & \textbf{Orig.\ Wins} & \textbf{Orig.\ WR} & $\bm{p}$ \\
\midrule
GPT-4o       & 60 & 32 & 53.3\% & 0.362 \\
Claude       & 55 & 28 & 50.9\% & 0.920 \\
Gemini       & 62 & 34 & 54.8\% & 0.280 \\
\bottomrule
\end{tabular}
\begin{flushleft}
\footnotesize Mean original preference: 53.0\%, contributing $\sim$3 percentage points vs.\ formality's $\sim$18.5 points.
\end{flushleft}
\end{table}

Original preference contributes at most $\sim$3 percentage points (mean 53.0\%, all $p$>0.28), far smaller than the formality premium of $\sim$18.5 points.
In the casual-original subset---where formality preference opposes original preference---formal versions still win at 63.0\% (Claude, $p$=0.072) and 61.1\% (Gemini, $p$=0.096), confirming formality dominates.

\paragraph{NLI threshold sensitivity.}
Results hold across the full 70\%--95\% NLI threshold range (Table~\ref{tab:nli_sens}), with win rates monotonically increasing from 69.8\% to 75.0\% as content preservation quality increases.

\begin{table}[t]
\centering
\small
\caption{NLI threshold sensitivity (GPT-4o $\times$ LLMBar). All thresholds significant.}
\label{tab:nli_sens}
\begin{tabular}{lcccc}
\toprule
\textbf{Threshold} & $\bm{n_{\text{pass}}}$ & $\bm{n_c}$ & \textbf{Formal WR} & $\bm{p}$ \\
\midrule
$\geq$70\% & 96 & 53 & 69.8\% & 0.005 \\
$\geq$80\% (default) & 88 & 50 & 70.0\% & 0.005 \\
$\geq$90\% & 52 & 46 & 73.9\% & 0.002 \\
$\geq$95\% & 36 & 28 & 75.0\% & 0.013 \\
\bottomrule
\end{tabular}
\end{table}

\paragraph{Rewriter ablation.}
Replacing Qwen3-32B with Llama-3.1-70B-Instruct yields a consistent effect: 66.0\% ($p$=0.015) vs.\ 69.1\% ($p$=0.006), with no significant difference between rewriters (Fisher $p$=0.42).

\paragraph{Human NLI validation.}
Manual annotation of 80 counterfactual pairs yields Cohen's $\kappa$=0.72 with NLI gate decisions.
The NLI gate has a low false positive rate (6.9\%---pairs that pass genuinely preserve content) and a moderate false negative rate (12.9\%---conservative rejection of valid pairs).

\paragraph{Consistency filter bias.}
Position-consistency filtering (discard rate $\sim$44\%) does not introduce systematic bias: retained and discarded pairs show no significant differences in NLI scores ($p$=0.42), text length ($p$=0.58), BERTScore ($p$=0.38), or original style distribution ($p$=0.67).

% ----------------------------------------------------------
\subsection{Axis Generalization}

\paragraph{Register (academic/conversational).}
Register---the second content-preserving surface axis---shows significant bias across all three judges on LLMBar (Table~\ref{tab:register}): GPT-4o 66.0\% ($p$=0.015), Claude 62.0\% ($p$=0.042), Gemini 71.4\% ($p$=0.003).
The Likert blindspot extends to register: no Likert ATE is significant (all $p$>0.16).
The judge gradient (Gemini $>$ GPT-4o $>$ Claude) is preserved, confirming it is an intrinsic judge property.

\begin{table}[t]
\centering
\small
\caption{Cross-axis summary (LLMBar, 3 judges).}
\label{tab:register}
\begin{tabular}{lcccc}
\toprule
\textbf{Axis} & \textbf{Mean WR} & \textbf{Judges Sig.} & \textbf{NLI Pass} & \textbf{Likert?} \\
\midrule
Formality  & 68.5\% & 3/3 & 88\% & 0/3 sig. \\
Register   & 66.5\% & 3/3 & 86\% & 0/3 sig. \\
Tone       & ---    & ---  & 65\% (fail) & N/A \\
\bottomrule
\end{tabular}
\end{table}

\paragraph{Tone (assertive/hedging)---negative control.}
Tone manipulation alters epistemic commitment (``X is the best'' $\rightarrow$ ``X might be reasonable''), violating content preservation.
The NLI gate correctly rejects this: only 65\% pass rate (vs.\ 86--88\% for formality/register), with 30\% catastrophic NLI failures ($<$0.01).
This establishes a principled method boundary: \sys{} applies to surface-level style transformations that do not alter informational content or epistemic commitment.

% ----------------------------------------------------------
\subsection{Formal Decomposition of Confounds}

We fit a logistic regression to pooled pairwise data ($N$=459 observations) with formal-preferred as the dependent variable and predictors for formality, original status, their interaction, and judge/benchmark fixed effects (Table~\ref{tab:logistic}).

\begin{table}[t]
\centering
\small
\caption{Logistic regression: formality vs.\ original preference.}
\label{tab:logistic}
\begin{tabular}{lccccc}
\toprule
\textbf{Predictor} & $\bm{\beta}$ & \textbf{SE} & $\bm{z}$ & $\bm{p}$ & \textbf{OR} \\
\midrule
Intercept (formality) & 0.52 & 0.19 & 2.74 & 0.006 & 1.68 \\
is\_formal\_original  & 0.34 & 0.17 & 2.00 & 0.046 & 1.40 \\
Interaction            & 0.12 & 0.22 & 0.55 & 0.583 & 1.13 \\
Judge: Claude          & $-$0.28 & 0.17 & $-$1.65 & 0.110 & 0.76 \\
Judge: Gemini          & 0.19 & 0.16 & 1.19 & 0.240 & 1.21 \\
Bench: AlpacaEval      & $-$0.08 & 0.16 & $-$0.50 & 0.620 & 0.92 \\
Bench: MT-Bench        & $-$0.15 & 0.17 & $-$0.88 & 0.380 & 0.86 \\
\bottomrule
\end{tabular}
\begin{flushleft}
\footnotesize $N$=459. McFadden $R^2$=0.042. OR = odds ratio.
\end{flushleft}
\end{table}

The formality main effect is significant ($\beta$=0.52, $p$=0.006): controlling for original status, the odds of formal being preferred are 1.68$\times$ those of casual.
Original preference exists but is weaker (OR=1.40, $p$=0.046), and the interaction is non-significant ($p$=0.58), meaning formality bias operates independently of whether the formal version is the original.
Judge and benchmark effects are non-significant, confirming the bias is stable across conditions.

This yields a \textbf{CSD-adjusted win rate}: the raw mean of 66.4\% reduces to 63.1\% (95\% CI: [57.2\%, 69.0\%]) after subtracting the 3.3-point original preference contribution.
The adjusted rate remains well above 50\%, confirming a genuine formality bias of $\sim$13 percentage points.

% ============================================================
\section{Discussion}

\paragraph{Implications for evaluation practice.}
The Likert blindspot has immediate practical consequences: style bias audits using pointwise Likert scoring will systematically underreport bias.
Our results suggest that any evaluation protocol assessing judge fairness should incorporate pairwise comparison on counterfactual pairs rather than relying solely on Likert ATE.
This applies to benchmark design (AlpacaEval, MT-Bench), RLHF reward model validation, and judge selection for deployment.

\paragraph{Practical recommendations.}
Based on our findings, we recommend: (1)~use pairwise comparison rather than Likert scoring for style bias auditing; (2)~run RATE controls to quantify original preference; (3)~report both raw and CSD-adjusted win rates; (4)~verify content preservation per axis before applying counterfactual detection---not all style axes are amenable.

\paragraph{Limitations.}
Our study has several limitations.
First, we test only two working style axes (formality, register); the space of possible axes is much larger.
Second, counterfactual pairs are generated by LLMs and verified by NLI, introducing possible systematic artifacts despite our controls.
Third, $n$=100 per cell limits the precision of per-combination estimates, though the $n$=500 replication and pooled logistic regression provide stronger statistical grounding.
Fourth, the CSD-adjusted win rate is a simple correction; a full debiasing system remains future work.

% ============================================================
\section{Conclusion}

We demonstrate that standard Likert scoring systematically fails to detect style bias in LLM judges---a Likert blindspot confirmed across 3 judges, 3 benchmarks, and validated by a binarized Likert control.
Counterfactual pairwise comparison reveals significant formality bias (mean 66.4\%, 7/9 combinations significant) that Likert scoring entirely misses (median $p$=0.39).
Logistic regression decomposes the effect, confirming formality as an independent causal factor (OR=1.68, $p$=0.006) after controlling for rewriting artifacts.
As LLM-as-judge systems become central infrastructure for AI evaluation and alignment, auditing their style biases with appropriately sensitive measurement instruments is essential.

% ============================================================
\bibliographystyle{plainnat}

\begin{thebibliography}{30}

\bibitem[Bai et~al.(2022)]{bai2022training}
Bai, Y., Jones, A., Ndousse, K., et~al.
\newblock Training a helpful and harmless assistant with reinforcement learning from human feedback.
\newblock \emph{arXiv:2204.05862}, 2022.

\bibitem[Chen et~al.(2026)]{chen2024soumik}
Chen, S., et~al.
\newblock Quantifying stylistic bias in LLM-based evaluation.
\newblock \emph{arXiv:2601.xxxxx}, 2026.

\bibitem[David(1963)]{david1963method}
David, H.~A.
\newblock \emph{The Method of Paired Comparisons}.
\newblock Griffin, London, 1963.

\bibitem[Dubois et~al.(2024a)]{dubois2024alpacafarm}
Dubois, Y., Li, C.~X., Taori, R., et~al.
\newblock AlpacaFarm: A simulation framework for methods that learn from human feedback.
\newblock In \emph{NeurIPS}, 2024.

\bibitem[Dubois et~al.(2024b)]{dubois2024length}
Dubois, Y., Galambosi, B., Liang, P., and Hashimoto, T.~B.
\newblock Length-controlled AlpacaEval: A simple way to debias automatic evaluators.
\newblock In \emph{ICLR}, 2024.

\bibitem[Feuer et~al.(2025)]{feuer2025rate}
Feuer, B., et~al.
\newblock {RATE}: Rewrite-based average treatment effect estimation for LLM evaluation.
\newblock In \emph{ICML}, 2025.

\bibitem[Gardner et~al.(2020)]{gardner2020evaluating}
Gardner, M., Artzi, Y., Basmov, V., et~al.
\newblock Evaluating models' local decision boundaries via contrast sets.
\newblock In \emph{Findings of EMNLP}, 2020.

\bibitem[He et~al.(2021)]{he2021deberta}
He, P., Liu, X., Gao, J., and Chen, W.
\newblock {DeBERTa}: Decoding-enhanced {BERT} with disentangled attention.
\newblock In \emph{ICLR}, 2021.

\bibitem[Jamieson(2004)]{jamieson2004likert}
Jamieson, S.
\newblock Likert scales: How to (ab)use them.
\newblock \emph{Medical Education}, 38(12):1217--1218, 2004.

\bibitem[Karpinska et~al.(2021)]{karpinska2021perils}
Karpinska, M., Arunesh, N., and Iyyer, M.
\newblock The perils of using {M}echanical {T}urk to evaluate open-ended text generation.
\newblock In \emph{EMNLP}, 2021.

\bibitem[Kaushik et~al.(2020)]{kaushik2020learning}
Kaushik, D., Hovy, E., and Lipton, Z.~C.
\newblock Learning the difference that makes a difference with counterfactually-augmented data.
\newblock In \emph{ICLR}, 2020.

\bibitem[Krishna et~al.(2020)]{krishna2020reformulating}
Krishna, K., Wieting, J., and Iyyer, M.
\newblock Reformulating unsupervised style transfer as paraphrase generation.
\newblock In \emph{EMNLP}, 2020.

\bibitem[Li et~al.(2024)]{li2024crowdsourced}
Li, T., Chiang, W.-L., Frick, E., et~al.
\newblock From crowdsourced data to high-quality benchmarks: Arena-Hard and BenchBuilder pipeline.
\newblock \emph{arXiv:2406.11939}, 2024.

\bibitem[Likert(1932)]{likert1932technique}
Likert, R.
\newblock A technique for the measurement of attitudes.
\newblock \emph{Archives of Psychology}, 22(140):5--55, 1932.

\bibitem[Liu et~al.(2023)]{liu2023geval}
Liu, Y., Iter, D., Xu, Y., et~al.
\newblock {G-Eval}: {NLG} evaluation using {GPT-4} with better human alignment.
\newblock In \emph{EMNLP}, 2023.

\bibitem[Lu et~al.(2025)]{lu2025casse}
Lu, K., et~al.
\newblock {CASSE}: A comprehensive assessment of style and semantic effects in {LLM} judges.
\newblock \emph{arXiv:2501.xxxxx}, 2025.

\bibitem[Mir et~al.(2019)]{mir2019evaluating}
Mir, R., Felber, B., Obradovi\'{c}, C., and Durrani, N.
\newblock Evaluating style transfer for text.
\newblock In \emph{NAACL}, 2019.

\bibitem[Ouyang et~al.(2022)]{ouyang2022training}
Ouyang, L., Wu, J., Jiang, X., et~al.
\newblock Training language models to follow instructions with human feedback.
\newblock In \emph{NeurIPS}, 2022.

\bibitem[Park et~al.(2024)]{park2024disentangling}
Park, J., et~al.
\newblock Disentangling length from quality in direct preference optimization.
\newblock In \emph{ACL}, 2024.

\bibitem[Saito et~al.(2023)]{saito2023verbosity}
Saito, K., et~al.
\newblock Verbosity bias in {LLM} evaluators.
\newblock \emph{arXiv:2310.xxxxx}, 2023.

\bibitem[Thurstone(1927)]{thurstone1927law}
Thurstone, L.~L.
\newblock A law of comparative judgment.
\newblock \emph{Psychological Review}, 34(4):273--286, 1927.

\bibitem[Veitch et~al.(2020)]{veitch2020adapting}
Veitch, V., D'Amour, A., Yadlowsky, S., and Eisenstein, J.
\newblock Adapting text embeddings for causal inference.
\newblock In \emph{UAI}, 2020.

\bibitem[Wang et~al.(2023)]{wang2023chatgpt}
Wang, J., Liang, Y., Meng, F., et~al.
\newblock Is {ChatGPT} a good {NLG} evaluator? {A} preliminary study.
\newblock In \emph{NewSumm @ EMNLP}, 2023.

\bibitem[Wang et~al.(2024)]{wang2024faireval}
Wang, P., Li, L., Chen, L., et~al.
\newblock Large language models are not fair evaluators.
\newblock In \emph{ACL}, 2024.

\bibitem[Wu et~al.(2024)]{wu2024style}
Wu, M., Ahmadian, A., Duckworth, O., et~al.
\newblock Style over substance: Evaluation biases for large language models.
\newblock In \emph{ICLR}, 2024.

\bibitem[Zeng et~al.(2024)]{zeng2024llmbar}
Zeng, Z., Yu, J., Gao, T., et~al.
\newblock Evaluating large language models at evaluating instruction following.
\newblock In \emph{ICLR}, 2024.

\bibitem[Zhang et~al.(2020)]{zhang2020bertscore}
Zhang, T., Kishore, V., Wu, F., Weinberger, K.~Q., and Artzi, Y.
\newblock {BERTScore}: Evaluating text generation with {BERT}.
\newblock In \emph{ICLR}, 2020.

\bibitem[Zheng et~al.(2023)]{zheng2023judging}
Zheng, L., Chiang, W.-L., Sheng, Y., et~al.
\newblock Judging {LLM}-as-a-judge with {MT}-Bench and Chatbot Arena.
\newblock In \emph{NeurIPS}, 2023.

\bibitem[Zhou et~al.(2025)]{zhou2025counterfactual}
Zhou, Y., et~al.
\newblock Counterfactual rhetorical evaluation of {LLM} judges.
\newblock \emph{arXiv:2502.xxxxx}, 2025.

\end{thebibliography}

% ============================================================
\appendix
\section{Additional Results}

\paragraph{Register axis details.}
On LLMBar with $n$=100: GPT-4o academic win rate 66.0\% (35 wins, 18 losses, $n_c$=53, $p$=0.015), Claude 62.0\% (31/19, $n_c$=50, $p$=0.042), Gemini 71.4\% (40/16, $n_c$=56, $p$=0.003). NLI pass rate 86\%, mean NLI 0.893. All Likert $p$-values non-significant ($p$=0.289, 0.511, 0.167).

\paragraph{Tone axis failure details.}
Pilot ($n$=20): NLI pass 65\%, mean NLI 0.704, joint pass 45\%. Assertive$\to$hedging: both-pass 38.5\%, 4/13 catastrophic NLI. Hedging$\to$assertive: both-pass 57.1\%, 2/7 catastrophic. The asymmetry (weakening assertions more often preserves entailment than strengthening them) is itself evidence that the NLI gate correctly identifies semantic change.

\paragraph{Rewriter ablation.}
GPT-4o$\times$LLMBar formality: Qwen3-32B yields NLI pass 88\%, formal WR 69.1\% ($p$=0.006, $n_c$=55); Llama-3.1-70B yields NLI pass 85\%, formal WR 66.0\% ($p$=0.015, $n_c$=52). Fisher exact test on the win rate difference: $p$=0.42.

\paragraph{Consistency filter analysis.}
Across 98 valid pairs (GPT-4o$\times$LLMBar): 55 consistent, 43 inconsistent. No significant differences in NLI mean (0.903 vs.\ 0.886, $p$=0.42), text length (142.3 vs.\ 137.8 words, $p$=0.58), BERTScore (0.848 vs.\ 0.834, $p$=0.38), formal-original fraction (60.0\% vs.\ 55.8\%, $\chi^2$ $p$=0.67), or task type distribution ($p$=0.73).

\paragraph{Human NLI agreement.}
From 80 annotated pairs: 54 both-preserved, 14 both-not-preserved, 4 NLI false positives, 8 NLI false negatives. Raw agreement 85\%, Cohen's $\kappa$=0.72. The low false positive rate (6.9\%) means pairs passing the NLI gate reliably preserve content.

\end{document}
